\section{CF-MoA: Sparse Orchestration of Five Functional Roles}
\label{sec:method}

CF-MoA implements five functional agent roles: adding support (A1), withdrawing support (A2), comparing candidates (A3), deriving consequences (A4), and preserving task-relevant judgments (A5). These roles correspond to R1--R5 and operate through a sparse controller. The adopted configuration, denoted B, selects an applicable operation from the current input's capabilities and returns a native answer (Figure~\ref{fig:architecture}). The base system supplies its current answer and actual inference context; agents share backbone weights and, where applicable, deterministic components.
% Evidence: docs/research/cf_moa.md, adopted_B_five_operations_v1;
% docs/paper/research_dossier.md, Section 5.9.

\subsection{Capability routing and native readout}

The controller selects operations by task format and available resources. Its fixed priority routes diagnosis-format inputs to A3, eligible rule tasks to the shared A1/A2 kernel, inputs providing a complete executable model to A4, and supported single-choice, robustness, or relation tasks with saved original messages to A5. Other inputs retain the base response. Routing uses visible options and question structure alongside format and resource availability; source identifiers, benchmark relation labels, gold answers, and paired endpoints are excluded. Capability-based routing allows a component developed on one functional subset to operate on compatible inputs elsewhere in the benchmark.

Each role operates within a defined domain: finite prescribing rules for A1/A2, a fixed diagnosis catalog for A3, a supplied executable mechanism and state for A4, and supported native decision formats for A5. In the complete v5 evaluation, each backbone routes 3,560 inputs to A3, 123 to A1/A2, 10,178 to A5, and 44 to passthrough. A4 has no qualifying clinical v5 input; its supplied-model evaluation is reported separately. The controller records the selected route, requests, operation traces, and complete returned response.
% Evidence: five_agent_framework_audit.json; complete v5 coverage in
% docs/research/cf_moa_experiment_report.md, Sections 16.6--16.13.

\subsection{A1/A2: scoped support recomputation}

A1 and A2 recompute support with a shared kernel of 71 manually encoded rules drawn from selected recommendations in the 2023 American Geriatrics Society Beers Criteria~\citep{ags2023beers}. Eligibility checks require a supported multiple-choice format, an age of at least 65 in the implemented question pattern, and a recognized medication option. The kernel identifies the requested action and scope, such as avoidance or renal dose adjustment. One call to the original backbone extracts the required patient facts and evidence sentence identifiers from the numbered current question. The extraction input excludes the base answer and retrieved documents. Missing or rejected evidence produces an unknown value.
% Evidence: docs/research/r2_support_head.md, execution guide Sections 2--3.

The executor evaluates conditions, exceptions, and independent support paths using \textsc{met}, \textsc{contradicted}, and \textsc{unknown}. Conjunction combines requirements within a path; disjunction preserves independent paths. Consequently, invalidating one risk does not remove an option supported by another satisfied risk. Let $S_0$ be the initial valid option set and let $M,C,U$ partition options by their aggregate status. A1 adds $M\setminus S_0$; A2 removes $C\cap S_0$. For ordinary valid sets, the result is
\begin{equation}
 S^{\star}=M\cup(S_0\cap U).
 \label{eq:support-update}
\end{equation}
The native assembler additionally handles exclusive ``None'' options, invalid initial answers, and empty outputs. Both roles therefore use one fact extraction and one execution, while retaining separate addition and withdrawal traces.
% Evidence: docs/research/cf_moa.md, minimal five-operation B;
% docs/research/r2_support_head.md, three-valued execution and assembly.

\subsection{A3: catalog-based candidate comparison}

A3 compares all 49 disease names in the public DDXPlus catalog~\cite{ddxplus}. The same backbone reads the current patient and the base system's actual context, with diagnoses bound to scoreable codes under six fixed catalog orders. For each order, it subtracts a content-free log-score prior from the patient-conditioned code score, following the general calibration principle of separating contextual evidence from answer preferences~\cite{zhao2021calibrate}. The prior is the logarithm of the mean of three content-free candidate distributions, computed separately for each order; its subtraction coefficient is fixed at one. A3 selects the highest corrected score per order and votes across orders, breaking vote ties using average corrected probabilities. The selected catalog entry is mapped back to the native diagnosis. The catalog supplies candidates for this decision stage; the scorer uses its separate, frozen normalization dictionary.
% Evidence: docs/research/r3_candidate_comparison.md, retained catalog/calibration;
% r3_catalog_calibration.py and r3_catalog_distribution_head.py.

\subsection{A4: executing a local intervention}

A4 compiles an intervention into a typed local patch and executes it to compute its consequences. Its input supplies the executable model, parameters, complete initial state, and factual action schedule. The patch preserves all unmodified actions. The numerical executor runs factual and patched trajectories from the same starting state and projects their consequences onto the requested time and outcome. Our implementation uses the public glucose--insulin simulator~\cite{simglucose}. Local patches reduce the need to regenerate a complete action plan, allowing intervention fidelity to be evaluated against a same-tool full-plan control. The returned consequence follows from the supplied mechanism and executed actions.
% Evidence: docs/research/r4_consequence_head.md, typed-patch delivery and
% frozen 24-profile validation; source simulator dependency is cited above.

\subsection{A5: task-adapted robust readout}

A5 implements task-adapted robust readout through component F. For ordinary supported tasks, F preserves the complete original messages and samples up to five native responses at temperature $0.7$, seeds 42--46, and a 2,048-token output limit. It votes over valid answer keys, stops once one key receives three votes, and returns the earliest complete response carrying a winning key. Final vote ties use earliest occurrence; if no vote is valid, the first response remains invalid. This adapts sampling and answer aggregation to the native output contract~\cite{wang2023selfconsistency}.

For explicit premise--hypothesis tasks with two polarity options, F uses a natural language inference (NLI) branch. It scores both A/B label mappings directly, generates one current-question analysis, and scores both mappings again conditioned on that analysis. Each two-label distribution is normalized before the four semantic distributions are averaged. The highest mean determines the answer; exact ties follow native option order. A final model call generates the complete native response with its answer field constrained to that decision. This path uses four one-token score callbacks and two complete generations at temperature zero, preserving the original evidence and required auxiliary fields. Missing code scores require a matching score request and remain visible in the cost ledger.
% Evidence: docs/research/r5_invariance_head.md, independent combination interface;
% Documents/Codex/2026-09-18/r5_head/delivery/r5_head.py, optimize/decide_paths.
